\section{Large Language Monkeys：重复采样的 Scaling Law}

\subsection{生成与验证的两阶段系统}

Repeated sampling 先从同一模型独立采样多个候选，再交给 verifier 选择。代码与形式化证明尤其适合这种方法，因为单元测试、编译器和 proof checker 能提供接近确定性的自动反馈。

\videofigure{figures/generate-verify.jpg}{多候选生成与验证构成 inference scaling 的基本流水线。}{00:01:18--00:02:22}

假设问题 $i$ 的单次采样成功概率为 $p_i$，独立采样 $k$ 次的 oracle coverage 为：

$$
C_i(k)=1-(1-p_i)^k.
$$

\begin{itemize}
    \item $p_i$：模型对问题 $i$ 的单次成功概率；
    \item $k$：采样数；
    \item $C_i(k)$：候选集合至少包含一个正确解的概率。
\end{itemize}

这说明每道题自身呈指数饱和：初始采样带来最大收益，之后边际增益逐步下降。

\subsection{小模型为何能通过采样追上强模型}

只要 $p_i>0$，增加 $k$ 就能提高 coverage。实验中，较小模型通过大量采样在数学、代码和形式化证明上显著缩小与强模型单次推理的差距。

\videofigure{figures/repeated-sampling.jpg}{重复采样在多个 reasoning benchmark 上显著放大模型能力。}{00:02:22--00:04:06}

\begin{knowledgebox}{能力被“放大”还是被“发现”}
Repeated sampling 并未改变参数，也没有让模型学会新算法。它做的是从模型已有的输出分布中，提高抽到低概率正确轨迹的机会。因此更准确的说法是\textbf{发现并利用潜在能力}，而不是创造全新能力。
\end{knowledgebox}

\subsection{从单题指数曲线到数据集平均幂律}

课程给出的经验拟合是 exponentiated power law：

$$
\bar C(k)\approx 1-\exp(-a k^b),\qquad a>0,\;0<b<1.
$$

\begin{itemize}
    \item $\bar C(k)$：数据集上的平均 coverage；
    \item $a$：与模型和任务总体难度相关的尺度参数；
    \item $b$：控制 scaling 速度的指数；
    \item $k$：推理采样预算。
\end{itemize}

\videofigure{figures/inference-scaling-law.jpg}{平均 coverage 随采样数呈可预测的 exponentiated power law。}{00:05:18--00:07:06}

单个问题是 $1-(1-p_i)^k$，为什么对问题平均后会出现幂律？关键在于 $p_i$ 的分布：数据集中存在大量模型偶尔才能解出的难题，形成靠近零的长尾。

\videofigure{figures/scaling-across-models.jpg}{从小模型到大模型都观察到 inference scaling。}{00:07:07--00:09:40}

\subsection{长尾是平均幂律的充分条件}

若问题成功概率在 $p\rightarrow0$ 附近具有密度

$$
f(p)\propto p^{\alpha-1},
$$

则数据集平均失败率近似为：

$$
1-\bar C(k)=\mathbb{E}_{p}[(1-p)^k]
\approx\int_{0}^{\infty} e^{-kp}p^{\alpha-1}\,dp
\propto k^{-\alpha}.
$$

\begin{itemize}
    \item $f(p)$：不同问题的单次成功概率分布；
    \item $\alpha$：靠近零点的长尾指数；
    \item $k^{-\alpha}$：随着采样预算增长而衰减的平均失败率。
\end{itemize}

\videofigure{figures/long-tail-theorem.jpg}{大量低成功概率问题使数据集平均 pass@$k$ 呈幂律扩展。}{00:09:40--00:11:42}

\begin{importantbox}{长尾决定 inference scaling 的持续时间}
如果所有题的 $p_i$ 都较高，少量采样后就饱和；如果大量题的 $p_i$ 极低但非零，收益会延续到很大的 $k$。因此推理 scaling 的形状不仅由模型决定，也由 benchmark 难度分布决定。
\end{importantbox}

\subsection{计算分配范式的变化}

当推理计算可预测地换取成功率时，系统不必把所有预算都固化在预训练中。对少量高价值请求，投入大量 test-time compute 可能比训练并部署更大模型更经济。

\videofigure{figures/compute-allocation.jpg}{从“训练大模型、单次推理”转向“模型适中、按请求投入推理计算”。}{00:11:42--00:12:38}

\subsection{本章小结}

Repeated sampling 的收益来自问题成功概率的长尾。它对模型规模具有广泛适用性，但它证明的是 oracle coverage 可扩展；若没有可靠 verifier，正确候选仍可能无法转化为最终输出。

